\setcounter{page}{136}
\setcounter{figure}{0}
\setcounter{chapter}{15}

\chapter{Momentum angular}
\label{cap:cap15}
\vspace{-10ex}

\raggedright

Hasta ahora hemos explorado cómo se describen la rotación y la aceleración angular de un cuerpo rígido bajo la acción de una torca. Siguiendo el paralelismo entre el movimiento lineal y el rotacional, es natural preguntarse: ¿cómo se transfiere energía en un sistema que rota? ¿Qué papel juegan el trabajo y la potencia cuando el movimiento no es lineal, sino angular?

En esta sección abordaremos dos ideas fundamentales que amplían nuestro entendimiento de los sistemas rotacionales:


\begin{itemize}
\item Por un lado, el \textbf{trabajo rotacional} se asocia con la acción de una torca que provoca una rotación, y se define de forma análoga al trabajo lineal, pero considerando la variable angular. La \textbf{potencia} en este contexto nos permitirá cuantificar la rapidez con la que se realiza trabajo en un sistema que gira.


\item Por otro lado, introduciremos el concepto de \textbf{momentum angular}, que es la cantidad de movimiento asociada a un objeto en rotación. Esta magnitud tiene un papel tan fundamental como la cantidad de movimiento lineal (momentum lineal), especialmente cuando se estudian sistemas donde las torcas externas son pequeñas o inexistentes, permitiendo el uso del principio de conservación del momentum angular.
\end{itemize}


Estos conceptos no solo son clave para comprender fenómenos físicos como el giro de una hélice o la rotación de un planeta, sino que también tienen aplicaciones prácticas en ingeniería, biomecánica, deportes y tecnologías que involucran sistemas en rotación, como turbinas, discos duros y giroscopios.
\section{Trabajo y potencia en el movimento de rotaci\'on}
El trabajo efectuado por una torca constante al generar un desplazamiento angular sobre un cuerpo est\'a dado por
\begin{equation}
    W=\tau_z \Delta \theta
\end{equation}

Si una torca efectúa trabajo sobre un cuerpo rígido que gira, la energía cinética experimenta un cambio en una cantidad igual a ese trabajo, es decir
\begin{equation}
    W_{\Mbox{total}}=\Delta K_{\Mbox{rot}}=\frac{1}{2}I\omega_{\Mbox{2}}^2-\frac{1}{2}I\omega_{\Mbox{1}}^2
\end{equation}
La potencia debida a una torca que actúa sobre un cuerpo rígido se relaciona con la velocidad angular del cuerpo con respecto
al eje de rotación mediante
\begin{equation}
    P=\tau_z \omega_z
\end{equation}

\section{Momento angular}

La \textit{cantidad de movimiento rotacional} de una part\'icula puntual de masa $m$ y momento lineal $\vec{p}$ se define como
\begin{equation}
\vec{L}=\vec{r}\times\vec{p}=\vec{r}\times m \vec{v}.
\end{equation}

 Definido de esta forma, el tambi\'en llamado \textit{momentum angular}, es la cantidad f\'isica que se ve alterada cuando un cuerpo (r\'igido o no r\'igido) gira bajo el efecto de una torca, es decir

\begin{equation}
\sum \vec{\bm{\tau}}=\der{\vec{L}}{t}.
\end{equation}

En el caso de un cuerpo r\'igido que gira alrededor de un eje de simetr\'ia con velocidad angular $\vec{\omega}$
 \begin{equation}
\vec{L}=I\vec{\omega},
\end{equation}
donde $I$ es el momento de inercia rotacional respecto al eje de simetr\'ia. \\

Para un cuerpo r\'igido, en donde el eje de giro permanece fijo, se cumple que 
\begin{equation}
L_z=I_z\omega_z.
\end{equation}
por lo que 
\begin{equation}
\sum \tau_z=\der{L_z}{t}=I_z\alpha_z.
\end{equation}

\section{Conservaci\'on del momento angular}

 Si la torca neta que act\'ua sobre un cuerpo r\'igido es nula, es decir, si 
 $$ \sum \vec{\bm{\tau}}=\vec{0}$$
 se cumple que $\vec{L}$ es constante (se conserva). Por lo tanto, en ausencia de torcas externas, se cumple para un cuerpo r\'igido que gira sin cambiar la direcci\'on de su eje de giro que 
\begin{equation}
I_{\Mbox{z,1}}\omega_{{\Mbox{z,1}}}=I_{\Mbox{z,2}}\omega_{{\Mbox{z,2}}}.
\end{equation}
Esta caracter\'istica es la que permite, por ejemplo, que una persona aumente su velocidad de giro al \textit{redistribuir} su masa en una "pirueta" de patinaje (ver Figura \ref{fig:patinador}).
\begin{figure}[!h]
\begin{center}
%\begin{tabular}{lr}
%\includegraphics[height=2in]{ejemplorot1}&
\includegraphics[width=3in]{planeamiento_didactico/figuras/patinador.png}
%\end{tabular}
\caption{\textit{Patinador girando}. $I_{\Mbox{z,1}}>I_{\Mbox{z,2}}$, por lo que $\omega_{{\Mbox{z,2}}}>\omega_{{\Mbox{z,1}}}$.}
\label{fig:patinador}
\textcolor{c4}{\footnotesize{Fuente: Sears, F. W., Zemansky, M. W., Young, H. D. y Freedman, R. A.
(2013). F\'isica Universitaria. Volumen 1. D\'ecimo tercera Edici\'on. M\'exico: Pearson Educaci\'on}}
\end{center}
\end{figure}